\section{Objectives and Test Scope}

\subsection{Purpose}
This report records the executed verification campaign for the device or
subsystem under test. It is intended to make the test purpose, boundaries,
success criteria, and reporting logic obvious to reviewers before they interpret
plots or tables.

\highlightbox{
Use this template when the main question is \emph{what was tested, how it was
tested, and whether each acceptance criterion passed}. It is stronger than a
generic lab report when you need explicit execution traceability.
}

\subsection{Device Under Test}
\begin{table}[H]
    \centering
    \small
    \begin{tabular}{l p{9.7cm}}
        \toprule
        \textbf{Field} & \textbf{Value / Guidance} \\
        \midrule
        Item Name & Name the board, assembly, firmware build, or subsystem revision. \\
        Serial / Build & Record the serial number, build lot, or assembly identifier. \\
        Test Objective & State the precise capability being verified. \\
        Campaign Owner & Name the engineer or team responsible for execution. \\
        Test Window & State the date range or lab session window for the activity. \\
        \bottomrule
    \end{tabular}
    \caption{Recommended device-under-test metadata}
    \label{tab:dut-metadata}
\end{table}

\subsection{Acceptance Criteria}
\begin{itemize}
    \item list measurable pass/fail criteria before pasting captured data;
    \item tie each criterion to a requirement, risk, or previous design concern;
    \item state any explicit exclusions to prevent scope creep during review;
    \item identify which results are screening-only and which are release-gating.
\end{itemize}

\notebox{
If different benches or fixtures were used during the campaign, declare that in
this chapter. Mixed environments are acceptable as long as they are explicit.
}
